\chapter{Calcolo Integrale ed Equazioni Differenziali}

\section{Integrazione Indefinita}
\subsection{Primitive e integrale indefinito}
\subsubsection{Definizione e proprietà di linearità}
\textbf{Definizione:} Siano $f, F: I \to \mathbb{R}$ due funzioni definite su un intervallo $I$. La funzione $F(x)$ si dice \textit{primitiva} di $f(x)$ se $F$ è derivabile in $I$ e risulta $F'(x) = f(x)$ per ogni $x \in I$.
Per una conseguenza del Teorema di Lagrange, se $F(x)$ è una primitiva di $f(x)$, lo sono anche tutte e sole le funzioni della forma $F(x) + c$, con $c \in \mathbb{R}$.
\textbf{Definizione:} L'insieme di tutte le primitive di una funzione $f(x)$ si definisce \textit{integrale indefinito} di $f(x)$ e si denota con il simbolo $\int f(x) dx = F(x) + c$.
\textbf{Proprietà di linearità:} L'operatore di integrazione indefinita è lineare:
\[ \int [\alpha f(x) + \beta g(x)] dx = \alpha \int f(x) dx + \beta \int g(x) dx \quad \forall \alpha, \beta \in \mathbb{R} \]

\subsubsection{Integrali immediati e quasi immediati}
Derivano direttamente dalle regole di derivazione lette al contrario. Ad esempio: $\int x^\alpha dx = \frac{x^{\alpha+1}}{\alpha+1} + c$ (per $\alpha \neq -1$), $\int \frac{1}{x} dx = \ln|x| + c$, $\int e^x dx = e^x + c$.
Gli integrali \textit{quasi immediati} sfruttano la regola di derivazione della funzione composta. La struttura è $\int f'(g(x))g'(x)dx = f(g(x)) + c$. 
\textbf{Esempio:} $\int \frac{f'(x)}{f(x)} dx = \ln|f(x)| + c$.

\subsection{Metodi di integrazione}
Quando un integrale non è riconducibile alle forme elementari, si applicano tecniche analitiche per trasformare la funzione integranda.

\subsubsection{Integrazione per scomposizione}
Sfruttando la linearità, si scompone la funzione integranda in una somma algebrica di funzioni i cui integrali siano noti. Si applica tipicamente a frazioni in cui il numeratore è una somma e il denominatore un monomio, oppure utilizzando identità trigonometriche.

\subsubsection{Integrazione per sostituzione}
Se l'integrale è nella forma $\int f(x) dx$, si può operare il cambio di variabile $x = g(t)$, con $g$ funzione derivabile e invertibile. Differenziando si ottiene $dx = g'(t) dt$, per cui:
\[ \int f(x) dx = \int f(g(t)) \cdot g'(t) dt \]
Risolto l'integrale in $t$, si opera la sostituzione inversa $t = g^{-1}(x)$ per tornare alla variabile originaria.

\subsubsection{Integrazione per parti}
Deriva dalla regola di derivazione del prodotto. Siano $f(x)$ e $g(x)$ due funzioni derivabili con derivata continua. Si ha:
\[ \int f(x)g'(x) dx = f(x)g(x) - \int f'(x)g(x) dx \]
La funzione $f(x)$ è detta \textit{fattore finito} (che verrà derivato), mentre $g'(x)dx$ è il \textit{fattore differenziale} (che verrà integrato per ricavare $g(x)$). La scelta opportuna di questi due fattori è cruciale per semplificare l'integrale risultante.

\subsubsection{Integrazione delle funzioni razionali fratte}
Si calcolano integrali del tipo $\int \frac{N(x)}{D(x)} dx$.
\begin{itemize}
    \item Se grado$(N) \ge$ grado$(D)$: Si esegue la divisione polinomiale $N(x) = D(x)Q(x) + R(x)$. L'integrale diventa $\int Q(x) dx + \int \frac{R(x)}{D(x)} dx$, abbassando il grado del numeratore.
    \item Se grado$(N) <$ grado$(D)$: Si applica il metodo di decomposizione in fratti semplici, fattorizzando il denominatore in $\mathbb{R}[x]$ (radici reali distinte, reali multiple, o complesse coniugate) e riscrivendo la frazione originaria come somma di frazioni parziali integrabili con logaritmi o arcotangenti.
\end{itemize}

\section{Integrazione Definita e Applicazioni}
\subsection{L'integrale definito}
\subsubsection{Problema dell'area e somme di Riemann}
Dato un intervallo $[a,b]$ e una funzione $f$ limitata e non negativa, si definisce una partizione di $[a,b]$ suddividendolo in $n$ sottointervalli. 
Siano $m_k$ e $M_k$ rispettivamente l'estremo inferiore e superiore di $f$ nel $k$-esimo sottointervallo di ampiezza $\Delta x_k$. Si costruiscono:
\begin{itemize}
    \item Somma integrale inferiore: $s_n = \sum_{k=1}^n m_k \Delta x_k$
    \item Somma integrale superiore: $S_n = \sum_{k=1}^n M_k \Delta x_k$
\end{itemize}
\textbf{Definizione (Integrabilità secondo Riemann):} Se l'estremo superiore delle somme inferiori coincide con l'estremo inferiore delle somme superiori al variare di tutte le possibili partizioni, tale valore comune è detto \textit{integrale definito} di $f$ da $a$ a $b$, indicato con $\int_a^b f(x) dx$. Geometricamente, rappresenta l'area del trapezoide sotteso dal grafico di $f$. Le funzioni continue sono sempre Riemann-integrabili.

\subsubsection{Definizione e proprietà (additività, linearità, monotonia)}
L'integrale definito associa a una funzione e a un intervallo un \textit{numero reale}.
\begin{itemize}
    \item \textbf{Additività rispetto all'intervallo:} Per ogni $c \in (a,b)$, $\int_a^b f(x) dx = \int_a^c f(x) dx + \int_c^b f(x) dx$.
    \item \textbf{Linearità:} $\int_a^b [\alpha f(x) + \beta g(x)] dx = \alpha \int_a^b f(x) dx + \beta \int_a^b g(x) dx$.
    \item \textbf{Monotonia (o Teorema del confronto):} Se $f(x) \le g(x)$ in $[a,b]$, allora $\int_a^b f(x) dx \le \int_a^b g(x) dx$.
\end{itemize}

\subsubsection{Teorema della media integrale}
Se $f(x)$ è continua in $[a,b]$, esiste almeno un punto $c \in (a,b)$ tale che:
\[ f(c) = \frac{1}{b-a} \int_a^b f(x) dx \]
Il valore a secondo membro è detto \textit{valor medio} della funzione nell'intervallo; geometricamente, $f(c)$ è l'altezza del rettangolo di base $(b-a)$ equivalente al trapezoide.

\subsection{Il Teorema Fondamentale del Calcolo Integrale}
\subsubsection{Funzione integrale e Teorema di Torricelli-Barrow}
Fissato un estremo $a$, si definisce \textit{funzione integrale} di $f(t)$ in $[a,b]$ la funzione $F(x) = \int_a^x f(t) dt$.
\textbf{Teorema Fondamentale (Torricelli-Barrow):} Se $f$ è continua in $[a,b]$, allora la funzione integrale $F(x)$ è derivabile in $(a,b)$ e risulta $F'(x) = f(x)$ per ogni $x \in (a,b)$. In sintesi, la funzione integrale è una primitiva della funzione integranda.

\subsubsection{Calcolo dell'integrale definito}
\textbf{Corollario (Formula fondamentale):} Sia $G(x)$ una qualsiasi primitiva di una funzione continua $f(x)$. Il calcolo pratico dell'integrale definito si esegue tramite l'identità:
\[ \int_a^b f(x) dx = \left[ G(x) \right]_a^b = G(b) - G(a) \]

\subsection{Applicazioni geometriche}
\subsubsection{Calcolo di aree di domini piani}
L'area della regione finita di piano compresa tra i grafici di due funzioni continue $f(x)$ e $g(x)$ (con $f(x) \ge g(x)$ per $x \in [a,b]$) e le rette verticali $x=a$ e $x=b$ è data da:
\[ \text{Area} = \int_a^b [f(x) - g(x)] dx \]
Se le curve si intersecano, l'integrale va spezzato o si integra il valore assoluto della differenza: $\int_a^b |f(x)-g(x)|dx$.

\subsubsection{Calcolo di volumi di solidi di rotazione (metodo dei dischi e dei gusci)}
\begin{itemize}
    \item \textbf{Rotazione attorno all'asse $x$ (Metodo delle sezioni/dischi):} Il volume del solido generato dalla rotazione completa del trapezoide individuato da $y=f(x)$ per $x \in [a,b]$ è:
    \[ V_x = \pi \int_a^b [f(x)]^2 dx \]
    \item \textbf{Rotazione attorno all'asse $y$ (Metodo dei gusci cilindrici):} Se il dominio piano compreso tra $f(x)$ e l'asse $x$ (con $0 \le a < b$) ruota attorno all'asse $y$:
    \[ V_y = 2\pi \int_a^b x \cdot |f(x)| dx \]
\end{itemize}

\subsubsection{Lunghezza di un arco di curva (cenni)}
Sia $f(x)$ una funzione dotata di derivata prima continua (funzione regolare) in $[a,b]$. La lunghezza dell'arco di curva da $x=a$ a $x=b$ è:
\[ L = \int_a^b \sqrt{1 + [f'(x)]^2} dx \]

\subsection{Integrazione impropria (Cenni)}
Estende il concetto di integrale definito a casi in cui l'intervallo è illimitato o la funzione integranda presenta asintoti verticali.
\subsubsection{Integrali estesi a intervalli illimitati}
Si sostituisce l'estremo infinito con una variabile $t$ e si passa al limite:
\[ \int_a^{+\infty} f(x) dx = \lim_{t \to +\infty} \int_a^t f(x) dx \]
Se il limite esiste finito, l'integrale improprio si dice \textit{convergente}; se è infinito, si dice \textit{divergente}; altrimenti è indeterminato.

\subsubsection{Integrali di funzioni illimitate}
Se $f(x)$ ha una singolarità nell'estremo $b$ (asintoto verticale), si pone:
\[ \int_a^b f(x) dx = \lim_{t \to b^-} \int_a^t f(x) dx \]

\section{Serie Numeriche e di Taylor}
\subsection{Serie numeriche}
\subsubsection{Definizione, convergenza, divergenza}
Data una successione $\{a_n\}$, si definisce successione delle \textit{somme parziali} $S_n = a_1 + a_2 + \dots + a_n = \sum_{k=1}^n a_k$. 
La \textbf{serie numerica} è il limite della successione $S_n$ al tendere di $n$ all'infinito: $\sum_{n=1}^\infty a_n = \lim_{n \to \infty} S_n$.
\begin{itemize}
    \item Se $\lim S_n = S \in \mathbb{R}$, la serie è \textbf{convergente} e la sua somma è $S$.
    \item Se $\lim S_n = \pm\infty$, la serie è \textbf{divergente}.
    \item Se il limite non esiste, la serie è \textbf{irregolare}.
\end{itemize}
\textbf{Condizione necessaria:} Affinché una serie converga, è necessario (ma non sufficiente) che $\lim_{n \to \infty} a_n = 0$.

\subsubsection{Serie geometrica e serie armonica}
\begin{itemize}
    \item \textbf{Serie geometrica} $\sum_{n=0}^\infty q^n$: Converge a $\frac{1}{1-q}$ se $|q| < 1$; diverge a $+\infty$ se $q \ge 1$; è irregolare se $q \le -1$.
    \item \textbf{Serie armonica generalizzata} $\sum_{n=1}^\infty \frac{1}{n^\alpha}$: Converge per $\alpha > 1$; diverge per $\alpha \le 1$. (Per $\alpha=1$ si ha la serie armonica semplice, storicamente nota per essere divergente pur avendo termine generale infinitesimo).
\end{itemize}

\subsubsection{Criteri di convergenza per serie a termini non negativi (Rapporto, Radice, Confronto)}
Data $\sum a_n$ con $a_n \ge 0$:
\begin{itemize}
    \item \textbf{Criterio del confronto asintotico:} Se $a_n \sim b_n$ per $n \to \infty$, allora le serie $\sum a_n$ e $\sum b_n$ hanno lo stesso carattere.
    \item \textbf{Criterio del rapporto (d'Alembert):} Se $\lim_{n \to \infty} \frac{a_{n+1}}{a_n} = l$, la serie converge per $l < 1$, diverge per $l > 1$.
    \item \textbf{Criterio della radice (Cauchy):} Se $\lim_{n \to \infty} \sqrt[n]{a_n} = l$, converge per $l < 1$, diverge per $l > 1$. 
    \item (Per $l=1$ i criteri del rapporto e della radice non forniscono informazioni).
\end{itemize}

\subsection{Sviluppi in serie (Cenni per il liceo)}
\subsubsection{Polinomi di Taylor e Maclaurin}
Ogni funzione $f$ dotata di $n$ derivate continue in un intorno di $x_0$ può essere approssimata localmente mediante il \textit{Polinomio di Taylor} di grado $n$:
\[ P_n(x) = f(x_0) + f'(x_0)(x - x_0) + \frac{f''(x_0)}{2!}(x - x_0)^2 + \dots + \frac{f^{(n)}(x_0)}{n!}(x - x_0)^n \]
L'errore commesso nell'approssimazione (Resto di Peano) è un infinitesimo di ordine superiore rispetto a $(x-x_0)^n$. Se $x_0 = 0$, si parla di polinomio e sviluppo di \textit{Maclaurin}.

\subsubsection{Sviluppi notevoli delle funzioni elementari}
Permettono di risolvere forme indeterminate complesse riducendole a rapporti polinomiali. Arrestati ai primi gradi, per $x \to 0$:
\begin{itemize}
    \item $e^x = 1 + x + \frac{x^2}{2!} + \dots + \frac{x^n}{n!} + o(x^n)$
    \item $\sin x = x - \frac{x^3}{3!} + \frac{x^5}{5!} + \dots + o(x^{2n+2})$ \quad (Solo potenze dispari)
    \item $\cos x = 1 - \frac{x^2}{2!} + \frac{x^4}{4!} + \dots + o(x^{2n+1})$ \quad (Solo potenze pari)
\end{itemize}

\section{Equazioni Differenziali}
\subsection{Generalità}
\subsubsection{Definizioni, ordine e soluzioni (generale, particolare)}
Un'equazione differenziale (ordinaria) è una relazione matematica che lega una variabile indipendente $x$, una funzione incognita $y(x)$ e alcune delle sue derivate $y', y'', \dots, y^{(n)}$.
\begin{itemize}
    \item \textbf{Ordine:} È il grado di derivazione massimo presente nell'equazione.
    \item \textbf{Integrale (soluzione) generale:} È la famiglia di funzioni che soddisfa l'equazione. Dipende da $n$ costanti arbitrarie (parametri).
    \item \textbf{Integrale particolare:} È una specifica soluzione ottenuta dall'integrale generale attribuendo valori particolari alle costanti.
\end{itemize}

\subsubsection{Il Problema di Cauchy}
Il Problema di Cauchy consiste nella ricerca di una e una sola soluzione particolare di un'equazione differenziale che soddisfi determinate \textit{condizioni iniziali}. Per un'equazione di ordine $n$, si assegnano i valori della funzione incognita e delle sue derivate fino all'ordine $(n-1)$ in un medesimo punto $x_0$. Sotto opportune ipotesi di regolarità (teorema di esistenza e unicità locale), il problema ammette soluzione unica.

\subsection{Equazioni differenziali del primo ordine}
La forma normale è $y' = f(x, y)$.

\subsubsection{Equazioni a variabili separabili}
Sono della forma $y' = f(x) \cdot g(y)$. Sfruttando la notazione differenziale di Leibniz $y' = \frac{dy}{dx}$, si separano le variabili:
\[ \frac{dy}{g(y)} = f(x) dx \implies \int \frac{1}{g(y)} dy = \int f(x) dx + c \]
Una volta integrati entrambi i membri, se possibile, si esplicita la $y$ in funzione della $x$.

\subsubsection{Equazioni differenziali lineari del primo ordine}
Hanno la forma $y' + a(x)y = b(x)$. Si risolvono individuando il \textit{fattore integrante} $e^{A(x)}$, dove $A(x) = \int a(x) dx$. Moltiplicando l'intera equazione per tale fattore, il primo membro diviene la derivata esatta del prodotto $y \cdot e^{A(x)}$. L'integrale generale si esprime mediante la formula chiusa:
\[ y(x) = e^{-A(x)} \left[ \int b(x) e^{A(x)} dx + c \right] \]

\subsection{Equazioni differenziali del secondo ordine}
Si considera il caso lineare: $ay'' + by' + cy = f(x)$, con $a, b, c$ coefficienti costanti reali ($a \neq 0$).

\subsubsection{Equazioni lineari a coefficienti costanti omogenee}
Se $f(x) = 0$, l'equazione è $ay'' + by' + cy = 0$. La forma dell'integrale generale dipende dal discriminante dell'equazione caratteristica algebrica associata $a\lambda^2 + b\lambda + c = 0$:
\begin{itemize}
    \item \textbf{$\Delta > 0$:} Due radici reali distinte $\lambda_1, \lambda_2$. L'integrale generale è $y(x) = c_1 e^{\lambda_1 x} + c_2 e^{\lambda_2 x}$.
    \item \textbf{$\Delta = 0$:} Radici reali coincidenti $\lambda_1 = \lambda_2 = \lambda$. L'integrale generale è $y(x) = c_1 e^{\lambda x} + c_2 x e^{\lambda x}$.
    \item \textbf{$\Delta < 0$:} Due radici complesse coniugate $\lambda = \alpha \pm i\beta$. Per la formula di Eulero, la base delle soluzioni reali conduce a $y(x) = e^{\alpha x} (c_1 \cos(\beta x) + c_2 \sin(\beta x))$.
\end{itemize}

\subsubsection{Equazioni lineari a coefficienti costanti non omogenee (cenni sui metodi risolutivi)}
La soluzione generale di un'equazione lineare completa è data dalla somma dell'integrale generale dell'equazione omogenea associata ($y_o$) e di un integrale particolare qualsiasi dell'equazione completa ($y_p$): $y(x) = y_o(x) + y_p(x)$.
La ricerca di $y_p(x)$ avviene per via algebrica tramite il \textit{metodo di somiglianza} (o dei coefficienti indeterminati). Se $f(x)$ ha forma polinomiale, esponenziale, sinusoidale o una loro composizione lineare, si imposta $y_p(x)$ con una forma analoga contenente parametri da determinare sostituendola nell'equazione di partenza. Un metodo più generale ma analiticamente gravoso è quello della \textit{variazione delle costanti arbitrarie}.

\section{Esercizi}

\subsection{Integrazione Indefinita: Immediati, Scomposizione e Fratti Semplici}
\begin{multicols}{2}
\begin{enumerate}
    % Integrali immediati e per scomposizione
    \item \tagprob{Int Immediato}{1} Calcola: $\int (3x^2 - 4x + 5) \, dx$
    \item \tagprob{Int Immediato}{1} Calcola: $\int \left(\frac{1}{x} + \frac{2}{x^2} - e^x\right) \, dx$
    \item \tagprob{Int Scomposizione}{1} Calcola: $\int \frac{x^2 - 3x + 1}{x} \, dx$
    \item \tagprob{Int Scomposizione}{1} Calcola: $\int \frac{x^3 + 8}{x+2} \, dx$ \textit{(Suggerimento: somma di cubi)}
    \item \tagprob{Int Quasi Immediato}{2} Calcola: $\int \sin(3x) \, dx$
    \item \tagprob{Int Quasi Immediato}{2} Calcola: $\int x e^{x^2} \, dx$
    \item \tagprob{Int Quasi Immediato}{2} Calcola: $\int \frac{2x}{x^2+4} \, dx$
    \item \tagprob{Int Quasi Immediato}{2} Calcola: $\int \tan x \, dx$ \textit{(Suggerimento: $\tan x = \frac{\sin x}{\cos x}$)}
    \item \tagprob{Int Scomposizione}{2} Calcola sfruttando identità trigonometriche: $\int \cos^2 x \, dx$
    \item \tagprob{Int Scomposizione}{3} Calcola: $\int \frac{1}{\sin^2 x \cos^2 x} \, dx$ \textit{(Suggerimento: $1 = \sin^2 x + \cos^2 x$)}
    
    % Integrazione Funzioni Razionali Fratte
    \item \tagprob{Int Fratti Sempl}{1} Divisione tra polinomi: $\int \frac{x^2+1}{x-1} \, dx$
    \item \tagprob{Int Fratti Sempl}{1} Denominatore radici reali distinte: $\int \frac{1}{x^2 - 4} \, dx$
    \item \tagprob{Int Fratti Sempl}{2} Denominatore radici reali distinte: $\int \frac{3x}{x^2 - 3x + 2} \, dx$
    \item \tagprob{Int Fratti Sempl}{2} Denominatore radice doppia: $\int \frac{x}{(x-1)^2} \, dx$
    \item \tagprob{Int Fratti Sempl}{2} Denominatore non fattorizzabile ($\Delta < 0$): $\int \frac{1}{x^2 + 4x + 5} \, dx$
    \item \tagprob{Int Fratti Sempl}{2} Denominatore con $\Delta < 0$ e numeratore grado 1: $\int \frac{2x+1}{x^2+9} \, dx$
    \item \tagprob{Int Fratti Sempl}{3} Numeratore grado superiore: $\int \frac{x^3 - 1}{x^2 - x - 2} \, dx$
    \item \tagprob{Int Fratti Sempl}{3} Metodo dei fratti complessi: $\int \frac{1}{x(x^2+1)} \, dx$
    \item \tagprob{Int Fratti Sempl}{3} Calcola: $\int \frac{2x^2+3}{(x^2+1)(x-2)} \, dx$
    \item \tagprob{Int Fratti Sempl}{3} Calcola (Hermite/Radici multiple): $\int \frac{1}{(x^2-1)^2} \, dx$
\end{enumerate}
\end{multicols}

\subsection{Integrazione per Sostituzione e per Parti}
\begin{multicols}{2}
\begin{enumerate}
    % Sostituzione
    \item \tagprob{Int Sostituzione}{1} Calcola ponendo $t = x^2$: $\int x \cos(x^2) \, dx$
    \item \tagprob{Int Sostituzione}{1} Calcola ponendo $t = \ln x$: $\int \frac{\ln x}{x} \, dx$
    \item \tagprob{Int Sostituzione}{2} Calcola ponendo $t = \sqrt{x-1}$: $\int \frac{x}{\sqrt{x-1}} \, dx$
    \item \tagprob{Int Sostituzione}{2} Calcola: $\int \frac{e^x}{e^{2x} + 1} \, dx$
    \item \tagprob{Int Sostituzione}{2} Calcola: $\int \frac{1}{x \sqrt{\ln x}} \, dx$
    \item \tagprob{Int Sostituzione}{3} Calcola ponendo $t = e^x$: $\int \frac{e^{2x}}{e^x - 1} \, dx$
    \item \tagprob{Int Sostituzione}{3} Calcola ponendo $x = \sin t$: $\int \sqrt{1 - x^2} \, dx$
    
    % Per Parti
    \item \tagprob{Int Parti}{1} Calcola scegliendo $f(x)=x$: $\int x e^x \, dx$
    \item \tagprob{Int Parti}{1} Calcola scegliendo $f(x)=x$: $\int x \sin x \, dx$
    \item \tagprob{Int Parti}{2} Calcola (fattore differenziale $1 \, dx$): $\int \ln x \, dx$
    \item \tagprob{Int Parti}{2} Calcola (fattore differenziale $1 \, dx$): $\int \arctan x \, dx$
    \item \tagprob{Int Parti}{2} Calcola applicando le parti due volte: $\int x^2 e^{-x} \, dx$
    \item \tagprob{Int Parti}{3} Calcola (integrale ciclico): $\int e^x \sin x \, dx$
    \item \tagprob{Int Parti}{3} Calcola: $\int x \ln(x^2) \, dx$
    \item \tagprob{Int Parti}{3} Calcola applicando prima sostituzione e poi parti: $\int \sin(\ln x) \, dx$
\end{enumerate}
\end{multicols}

\subsection{Integrazione Definita, Aree e Volumi}
\begin{multicols}{2}
\begin{enumerate}
    % Integrali Definiti di base
    \item \tagprob{Int Definito}{1} Calcola: $\int_1^2 (3x^2 - 1) \, dx$
    \item \tagprob{Int Definito}{1} Calcola: $\int_0^{\pi/2} \cos x \, dx$
    \item \tagprob{Int Definito}{2} Calcola applicando la sostituzione e aggiornando gli estremi: $\int_0^1 x e^{x^2} \, dx$
    \item \tagprob{Int Definito}{2} Calcola integrando per parti: $\int_1^e \ln x \, dx$
    \item \tagprob{Int Definito}{2} Calcola gestendo il valore assoluto: $\int_{-1}^2 |x| \, dx$
    
    % Calcolo di Aree
    \item \tagprob{Area Piana}{1} Calcola l'area della regione sottesa dalla parabola $y=4-x^2$ e l'asse $x$.
    \item \tagprob{Area Piana}{2} Calcola l'area compresa tra $y = x^2$ e la retta $y = x$.
    \item \tagprob{Area Piana}{2} Calcola l'area compresa tra $y = \sin x$ e l'asse $x$ nell'intervallo $[0, 2\pi]$.
    \item \tagprob{Area Piana}{3} Calcola l'area della regione finita compresa tra $y = e^x$, $y = e^{-x}$ e la retta $x = \ln 2$.
    \item \tagprob{Area Piana}{3} Calcola l'area compresa tra $y = \frac{1}{x}$ e la retta passante per i punti $(1,1)$ e $(2, 1/2)$.
    
    % Volumi e Teorema Media
    \item \tagprob{Volume Rotazione}{2} Calcola il volume del cono ottenuto ruotando $y = x$ per $x \in [0, h]$ attorno all'asse $x$.
    \item \tagprob{Volume Rotazione}{2} Calcola il volume del solido (paraboloide) ottenuto ruotando $y = \sqrt{x}$ per $x \in [0, 4]$ attorno all'asse $x$.
    \item \tagprob{Volume Rotazione}{3} Sfruttando il metodo dei gusci, calcola il volume ottenuto ruotando $y = 4-x^2$ (con $x \in [0,2]$) attorno all'asse $y$.
    \item \tagprob{Teorema Media}{2} Determina il valore medio della funzione $f(x) = x^2$ in $[0,3]$ e trova il punto $c$ garantito dal teorema.
    \item \tagprob{Funzione Integrale}{3} Data $F(x) = \int_0^x e^{t^2} dt$, scrivi l'equazione della retta tangente al grafico di $F(x)$ in $x=0$.
\end{enumerate}
\end{multicols}

\subsection{Serie Numeriche, Sviluppi di Taylor e Integrali Impropri}
\begin{multicols}{2}
\begin{enumerate}
    % Integrali Impropri
    \item \tagprob{Int Improprio}{1} Stabilisci se converge e, in caso, calcola: $\int_1^{+\infty} \frac{1}{x^2} \, dx$
    \item \tagprob{Int Improprio}{1} Mostra tramite calcolo diretto che $\int_1^{+\infty} \frac{1}{x} \, dx$ diverge.
    \item \tagprob{Int Improprio}{2} Stabilisci se converge e calcola: $\int_0^1 \frac{1}{\sqrt{x}} \, dx$
    \item \tagprob{Int Improprio}{3} Stabilisci se l'integrale \textit{campana di Gauss} (solo $x \ge 0$) converge: $\int_0^{+\infty} x e^{-x^2} \, dx$.
    
    % Serie Numeriche e Taylor
    \item \tagprob{Serie Geometrica}{1} Calcola la somma della serie $\sum_{n=0}^\infty \left(\frac{1}{3}\right)^n$.
    \item \tagprob{Serie Armonica}{1} Determina il carattere della serie $\sum_{n=1}^\infty \frac{1}{\sqrt{n}}$.
    \item \tagprob{Serie Criterio}{2} Usa il criterio del rapporto per studiare $\sum_{n=1}^\infty \frac{n}{2^n}$.
    \item \tagprob{Serie Criterio}{2} Usa il criterio della radice per studiare $\sum_{n=1}^\infty \left(\frac{n}{3n+1}\right)^n$.
    \item \tagprob{Serie Confronto}{3} Studia il carattere della serie usando il confronto asintotico: $\sum_{n=1}^\infty \sin\left(\frac{1}{n^2}\right)$.
    \item \tagprob{Polinomio Taylor}{2} Scrivi il polinomio di Maclaurin di grado 3 per $f(x) = e^{2x}$.
    \item \tagprob{Polinomio Taylor}{3} Usando gli sviluppi notevoli, calcola $\lim_{x \to 0} \frac{x - \sin x}{x^3}$.
\end{enumerate}
\end{multicols}

\subsection{Equazioni Differenziali}
\begin{multicols}{2}
\begin{enumerate}
    % A variabili separabili
    \item \tagprob{Eq Diff Separabili}{1} Risolvi determinando l'integrale generale: $y' = 2x y$
    \item \tagprob{Eq Diff Separabili}{1} Risolvi l'equazione differenziale: $y' = x^2 e^{-y}$
    \item \tagprob{Cauchy Separabili}{2} Risolvi il problema di Cauchy: $\begin{cases} y' = -y \\ y(0) = 5 \end{cases}$
    \item \tagprob{Cauchy Separabili}{2} Risolvi il problema di Cauchy: $\begin{cases} y' = y \cos x \\ y(\pi/2) = 1 \end{cases}$
    \item \tagprob{Eq Diff Separabili}{3} Trova la soluzione in forma implicita di $y' = \frac{x}{y}$.
    
    % Lineari del 1° ordine
    \item \tagprob{Eq Diff 1 Ordine}{2} Determina l'integrale generale: $y' + 2y = e^{-x}$
    \item \tagprob{Eq Diff 1 Ordine}{2} Risolvi: $y' - \frac{1}{x} y = x^2$
    \item \tagprob{Cauchy 1 Ordine}{3} Risolvi il problema di Cauchy: $\begin{cases} y' + y \tan x = \sin x \\ y(0) = 2 \end{cases}$ \textit{(per $x \in (-\pi/2, \pi/2)$)}
    
    % Lineari del 2° ordine a coefficienti costanti
    \item \tagprob{Eq Diff 2 Ordine}{1} Omogenea ($\Delta > 0$): $y'' - 5y' + 6y = 0$
    \item \tagprob{Eq Diff 2 Ordine}{1} Omogenea ($\Delta = 0$): $y'' - 4y' + 4y = 0$
    \item \tagprob{Eq Diff 2 Ordine}{2} Omogenea ($\Delta < 0$): $y'' + 9y = 0$ (Eq. dell'oscillatore armonico)
    \item \tagprob{Cauchy 2 Ordine}{2} Risolvi il problema di Cauchy: $\begin{cases} y'' - y = 0 \\ y(0) = 1 \\ y'(0) = 0 \end{cases}$
    \item \tagprob{Eq Diff Completa}{3} Usa il metodo di somiglianza per risolvere: $y'' - y' = 2x$
    \item \tagprob{Eq Diff Completa}{3} Usa il metodo di somiglianza per risolvere: $y'' + y = e^x$
    \item \tagprob{Eq Diff Mista}{3} Determina la famiglia di curve tali che la derivata seconda è costantemente uguale a 4 e che passano per l'origine.
\end{enumerate}
\end{multicols}